%
% sentinel-reap.tex
%
% Z Specification for Biff's Sentinel/Reaper Logout Mechanism
%
% Models sentinel-file writing, reaping, the own-key skip, and the
% liveness relationship between a sentinel and the KV row it refers to.
%
% Type-check: fuzz sentinel-reap.tex
%

\documentclass[a4paper,10pt,fleqn]{article}
\usepackage[margin=1in]{geometry}
\usepackage{fuzz}
\usepackage[colorlinks=true,linkcolor=blue,citecolor=blue,urlcolor=blue]{hyperref}

\begin{document}

\title{Biff Sentinel/Reaper Logout: A Z Specification}
\author{Punt Labs}
\date{August 2026}
\hypersetup{
  pdftitle={Biff Sentinel/Reaper Logout: A Z Specification},
  pdfauthor={Punt Labs},
  pdfsubject={Z Specification},
  pdfcreator={fuzz/probcli}
}
\maketitle

\tableofcontents
\newpage

%%----------------------------------------------------------------------
\section{Introduction}
\label{sec:intro}
%%----------------------------------------------------------------------

Every biff session owns exactly one row in the sessions KV bucket,
modelled in \texttt{nats-relay.tex} as $SessionStore$
(\texttt{sessions~:~SESSIONKEY~\pfun~KVVALUE}).  That specification is
deliberately silent on how a row is \emph{removed} when its owning
process goes away: $DeleteSession$'s commentary reads only ``triggered
by explicit delete or KV TTL expiry,'' and \texttt{session-model.tex}
lists ``heartbeats, TTLs, orphan detection, crash recovery'' among what
it does not cover.  The mechanism that fills that gap is the sentinel
and its reaper, implemented in \texttt{src/biff/server/app.py} and,
until this document, present nowhere in the specification corpus.

A sentinel is a file at
\texttt{\textasciitilde/.punt-labs/biff/sentinels/\{repo\}/\{session\_key\}}
(\texttt{\_write\_sentinel}), written just before a session's owning
process attempts to tear itself down.  Any running server's periodic
reaper (\texttt{\_reap\_loop}, every two seconds) or the initial call in
\texttt{\_active\_lifespan} at startup scans the sentinel directory
(\texttt{\_reap\_sentinels}) and, for each sentinel whose session key is
not its own, deletes the KV row, releases the TTY name, appends a
logout event, and removes the sentinel (\texttt{\_reap\_dead\_session}).
For a sentinel matching \emph{its own} session key, it discards the
sentinel without touching the row --- the ``own-key skip'':

\begin{quote}
\texttt{if session\_key in live\_keys:\\
\hphantom{xx}\# Our own prior incarnation --- consume the sentinel\\
\hphantom{xx}\# without a logout, reservation release, or KV delete.\\
\hphantom{xx}sentinel.unlink(missing\_ok=True)\\
\hphantom{xx}continue}
\end{quote}

Until the PR that motivates this specification (biff-7xd, \#350), only
the signal handler (\texttt{SIGTERM}/\texttt{SIGINT}/\texttt{SIGHUP})
wrote a sentinel.  A normal shutdown --- an MCP client disconnecting,
which runs \texttt{\_lifespan\_cleanup} through ordinary coroutine
control flow rather than a signal --- had no sentinel and so no reaper
coverage: a \texttt{delete\_session} call that timed out under the new
per-step bound (\texttt{\_TEARDOWN\_STEP\_TIMEOUT}, biff-7xd) silently
orphaned the row.  \#350 closes that gap by writing the same sentinel
unconditionally at the top of \texttt{\_lifespan\_cleanup}
(\texttt{\_write\_reap\_fallback\_sentinels}), \emph{before} the bounded,
best-effort teardown runs, so the sentinel mechanism is now the primary
logout path for both routes: signal-triggered and ordinary disconnect.

This specification models that mechanism from the point of view of one
distinguished, currently-running incarnation, and asks the question the
bead (biff-lx2) poses directly: is there a reachable state in which a
KV row has no live process behind it and no sentinel left to reclaim
it?  \S\ref{sec:modelcheck} answers in two parts.  The trivial part ---
a process killed with \texttt{SIGKILL} never gets to write a sentinel
at all --- is a known, already-documented gap (biff-teh's discovery
notes: ``SIGKILL bypasses the sentinel write, so force-killed sessions
linger in KV until the 120s liveness window expires''), backstopped by
KV TTL and the presence-liveness filter, and outside what a sentinel
mechanism can promise; the model reproduces it in one step and moves
on, following \texttt{nats-relay.tex}'s own precedent of leaving TTL
expiry unmodelled as a live transition.  The non-trivial part is new:
\S\ref{sec:race} exhibits a five-step counterexample, internal to the
mechanism \#350 ships, in which the reaping incarnation's \emph{own}
background reap loop races its \emph{own} fallback-sentinel write
during an ordinary, non-crashing shutdown --- no signal, no crash, and
no resume involved.

%%----------------------------------------------------------------------
\section{Basic Types}
%%----------------------------------------------------------------------

\begin{zed}
  [SESSIONKEY]
\end{zed}

$SESSIONKEY$ is the composite $\{user\}:\{tty\}$ identifier, exactly as
in \texttt{nats-relay.tex} --- redeclared here because each Z document
in this corpus is independently type-checked and so restates the basic
types it needs, not because the concept differs.  As in
\texttt{session-model.tex}'s $SESSION$, a $SESSIONKEY$ is stable across
\texttt{claude --resume}: the routing token (\texttt{tty\_hex}) is
derived from the Claude \texttt{session\_id}, so a resumed session
recovers the \emph{same} key rather than minting a fresh one.  That
stability is exactly what makes the own-key skip meaningful --- and, as
\S\ref{sec:race} shows, exactly what makes it unable to distinguish a
legitimately resumed incarnation from the very incarnation that wrote
the sentinel still running.

\begin{zed}
  ZBOOL ::= ztrue | zfalse
\end{zed}

%%----------------------------------------------------------------------
\section{Global Constants}
\label{sec:own}
%%----------------------------------------------------------------------

\begin{axdef}
  own : SESSIONKEY
\end{axdef}

$own$ is the session key of the incarnation whose perspective this
model takes --- the value of \texttt{state.session\_key} inside one
running \texttt{biff mcp} process.  It is a fixed, arbitrary element of
$SESSIONKEY$: the model does not care \emph{which} key it is, only that
every operation below agrees on it, exactly as \texttt{\_reap\_sentinels}
compares every sentinel it finds against the one value
\texttt{state.session\_key} for the life of the process.  Companion
sessions (\texttt{state.companion\_session\_key}) are a second,
independent instance of the same $own$-relative machinery --- the code
folds both into one \texttt{live\_keys} set --- and are not modelled
separately; every property proved of $own$ applies unchanged to a
companion key.

%%----------------------------------------------------------------------
\section{State}
%%----------------------------------------------------------------------

\begin{schema}{SentinelWorld}
  rowPresent : \power SESSIONKEY \\
  sentinelPresent : \power SESSIONKEY \\
  sentinelWriteAttempted : ZBOOL \\
  shutdownSignaled : ZBOOL \\
  releaseAttempted : ZBOOL \\
  exited : ZBOOL \\
  crashed : ZBOOL
\end{schema}

$rowPresent$ is the set of session keys with a live KV row --- the
domain of \texttt{nats-relay.tex}'s $SessionStore.sessions$, projected
down to presence, since this specification is not concerned with a
row's contents.  $sentinelPresent$ is the set of session keys with a
sentinel file currently on disk for this repo.  $sentinelWriteAttempted$
is a per-incarnation latch: \emph{this} incarnation of $own$ has called
\texttt{\_write\_reap\_fallback\_sentinels} or the signal handler's
sentinel write at least once.  It is a ghost, not an observable of the
real system, needed because \texttt{\_shutdown\_tasks} is called
\emph{after} the sentinel write in program order
(\texttt{\_lifespan\_cleanup}'s sequence is: write sentinel, append
logout events, \emph{then} shut down background tasks) regardless of
whether the sentinel it wrote is still sitting on disk by the time
$shutdownSignaled$ is set --- and a state predicate on
$sentinelPresent$ cannot express a fact about program order once the
file itself may already have been raced away.  $shutdownSignaled$
records whether $own$'s background \texttt{\_reap\_loop} has been told
to stop (its \texttt{asyncio.Event} set) --- once true, that task can no
longer tick and so can no longer run $OwnOrphanReapTick$ below.
$releaseAttempted$ is a second per-incarnation latch: $own$'s
\texttt{\_release\_session} has been called at least once, whether or
not \texttt{delete\_session} succeeded.  \texttt{\_lifespan\_cleanup}
always calls \texttt{\_release\_relay} --- and so always attempts
\texttt{\_release\_session} on $own$ --- between $shutdownSignaled$
becoming true and the function returning, so $ExitProcess$ below is
gated on it: a process cannot exit through this, the graceful path,
without release having at least been tried.  $exited$
records whether $own$'s OS process has actually terminated.  $crashed$
is a latch, in the manner of \texttt{session-model.tex}'s $lostF$ and
$misF$, flipped only by the one operation that models an unannounced
death (\S\ref{sec:crash}); it exists purely to let
\S\ref{sec:modelcheck} separate the trivial crash counterexample from
the non-trivial race one.

No structural invariant is stated in $SentinelWorld$ itself.  Unlike
\texttt{nats-relay.tex}'s $Connection$ invariant (which holds in every
reachable state), the natural candidate invariant here ---
``$exited = ztrue$ implies $shutdownSignaled = ztrue$'' --- is
\emph{false} of the system as a whole: that is precisely what a crash
is.  Following \texttt{session-model.tex}'s treatment of delivery
soundness, the properties this specification proves are stated as
reachability goals over $SentinelWorld$ (\S\ref{sec:modelcheck}), not as
a state invariant that would have to be weakened to accommodate a case
the mechanism was never meant to cover.

%%----------------------------------------------------------------------
\section{Initial State}
%%----------------------------------------------------------------------

The model starts with $own$ freshly registered: a live row, no
sentinel, no shutdown in progress, not exited, nothing crashed.

\begin{schema}{Init}
  SentinelWorld'
  \where
  rowPresent' = \{ own \} \\
  sentinelPresent' = \emptyset \\
  sentinelWriteAttempted' = zfalse \\
  shutdownSignaled' = zfalse \\
  releaseAttempted' = zfalse \\
  exited' = zfalse \\
  crashed' = zfalse
\end{schema}

%%----------------------------------------------------------------------
\section{Operations on $own$}
%%----------------------------------------------------------------------

%%----------------------------------------------------------------------
\subsection{WriteOwnSentinel}
%%----------------------------------------------------------------------

Corresponds to \texttt{\_write\_reap\_fallback\_sentinels} (called at
the top of \texttt{\_lifespan\_cleanup}) and to the sentinel write in
\texttt{\_active\_lifespan}'s \texttt{\_signal\_handler}: both call
\texttt{\_write\_sentinel(repo, state.session\_key)} and are, at this
level of abstraction, the same event.  It can fire more than once for
one incarnation --- a signal can arrive and then the MCP transport can
also detect disconnect --- which is why the effect is idempotent
($\cup$, not requiring absence first) and why $sentinelWriteAttempted$,
once set, is never cleared by this operation.

\begin{schema}{WriteOwnSentinel}
  \Delta SentinelWorld
  \where
  exited = zfalse \\
  %%
  sentinelPresent' = sentinelPresent \cup \{ own \} \\
  sentinelWriteAttempted' = ztrue \\
  %%
  rowPresent' = rowPresent \\
  shutdownSignaled' = shutdownSignaled \\
  releaseAttempted' = releaseAttempted \\
  exited' = exited \\
  crashed' = crashed
\end{schema}

%%----------------------------------------------------------------------
\subsection{OwnOrphanReapTick}
%%----------------------------------------------------------------------

The own-key skip.  Models one tick of $own$'s background
\texttt{\_reap\_loop} (or the explicit call in
\texttt{\_active\_lifespan} before that task exists) finding a sentinel
that names $own$ itself, and discarding it \emph{without} deleting the
row, releasing the TTY reservation, or appending a logout event ---
exactly the \texttt{if session\_key in live\_keys} branch quoted in
\S\ref{sec:intro}.  Enabled only while the loop has not yet been told to
stop ($shutdownSignaled = zfalse$), mirroring
\texttt{while not shutdown.is\_set()}.

\begin{schema}{OwnOrphanReapTick}
  \Delta SentinelWorld
  \where
  shutdownSignaled = zfalse \\
  own \in sentinelPresent \\
  %%
  sentinelPresent' = sentinelPresent \setminus \{ own \} \\
  %%
  rowPresent' = rowPresent \\
  sentinelWriteAttempted' = sentinelWriteAttempted \\
  shutdownSignaled' = shutdownSignaled \\
  releaseAttempted' = releaseAttempted \\
  exited' = exited \\
  crashed' = crashed
\end{schema}

%%----------------------------------------------------------------------
\subsection{SignalShutdown}
%%----------------------------------------------------------------------

Corresponds to \texttt{\_shutdown\_tasks} setting the
\texttt{asyncio.Event} that stops the poller, reaper, heartbeat, and KV
watcher tasks.  In \texttt{\_lifespan\_cleanup} this always happens
\emph{after} the fallback sentinel write and the logout-event appends,
which is why the precondition is $sentinelWriteAttempted = ztrue$ ---
the program-order fact --- rather than $own \in sentinelPresent$, which
by the time this fires may already have been raced away by
$OwnOrphanReapTick$.

\begin{schema}{SignalShutdown}
  \Delta SentinelWorld
  \where
  shutdownSignaled = zfalse \\
  sentinelWriteAttempted = ztrue \\
  %%
  shutdownSignaled' = ztrue \\
  %%
  rowPresent' = rowPresent \\
  sentinelPresent' = sentinelPresent \\
  sentinelWriteAttempted' = sentinelWriteAttempted \\
  releaseAttempted' = releaseAttempted \\
  exited' = exited \\
  crashed' = crashed
\end{schema}

%%----------------------------------------------------------------------
\subsection{ReleaseOwnSuccess}
%%----------------------------------------------------------------------

Corresponds to \texttt{\_release\_session} succeeding: \texttt{delete\_
session} returns without raising, so \texttt{\_remove\_sentinel} is
called unconditionally afterwards.  Removing an absent sentinel is a
no-op (\texttt{missing\_ok=True}), so the effect on $sentinelPresent$
is safe whether or not $OwnOrphanReapTick$ has already consumed it.

\begin{schema}{ReleaseOwnSuccess}
  \Delta SentinelWorld
  \where
  shutdownSignaled = ztrue \\
  own \in rowPresent \\
  %%
  rowPresent' = rowPresent \setminus \{ own \} \\
  sentinelPresent' = sentinelPresent \setminus \{ own \} \\
  releaseAttempted' = ztrue \\
  %%
  sentinelWriteAttempted' = sentinelWriteAttempted \\
  shutdownSignaled' = shutdownSignaled \\
  exited' = exited \\
  crashed' = crashed
\end{schema}

%%----------------------------------------------------------------------
\subsection{ReleaseOwnFails}
%%----------------------------------------------------------------------

Corresponds to \texttt{delete\_session} raising or timing out inside
\texttt{\_release\_session}'s \texttt{\_TEARDOWN\_STEP\_TIMEOUT} bound
(biff-7xd) --- a ``log and move on'' failure.  Per the code,
\texttt{\_remove\_sentinel} is only reached in the \texttt{else} branch
of the \texttt{try}, so on failure the sentinel --- if $OwnOrphanReapTick$
has not already removed it --- is left exactly as it was.

\begin{schema}{ReleaseOwnFails}
  \Delta SentinelWorld
  \where
  shutdownSignaled = ztrue \\
  own \in rowPresent \\
  %%
  rowPresent' = rowPresent \\
  sentinelPresent' = sentinelPresent \\
  releaseAttempted' = ztrue \\
  %%
  sentinelWriteAttempted' = sentinelWriteAttempted \\
  shutdownSignaled' = shutdownSignaled \\
  exited' = exited \\
  crashed' = crashed
\end{schema}

%%----------------------------------------------------------------------
\subsection{ExitProcess}
%%----------------------------------------------------------------------

The OS process backing $own$ actually terminates --- the end of
\texttt{\_lifespan\_cleanup}'s \texttt{finally} block, or (post
biff-teh) the kernel-level re-delivery of the caught signal.  Gated on
$shutdownSignaled = ztrue \land releaseAttempted = ztrue$:
\texttt{\_shutdown\_tasks} is awaited, and \texttt{\_release\_relay} ---
which always attempts \texttt{\_release\_session} on $own$ --- runs,
strictly before \texttt{\_lifespan\_cleanup} returns and the process can
exit through this, the graceful path.

\begin{schema}{ExitProcess}
  \Delta SentinelWorld
  \where
  shutdownSignaled = ztrue \\
  releaseAttempted = ztrue \\
  exited = zfalse \\
  %%
  exited' = ztrue \\
  %%
  rowPresent' = rowPresent \\
  sentinelPresent' = sentinelPresent \\
  sentinelWriteAttempted' = sentinelWriteAttempted \\
  shutdownSignaled' = shutdownSignaled \\
  releaseAttempted' = releaseAttempted \\
  crashed' = crashed
\end{schema}

%%----------------------------------------------------------------------
\subsection{CrashNoSentinel}
\label{sec:crash}
%%----------------------------------------------------------------------

\texttt{SIGKILL}, OOM, or power loss: the process is gone with no
warning and so with no chance to run \texttt{\_write\_sentinel} or
\texttt{\_shutdown\_tasks}.  This is the case biff-teh's own discovery
notes name explicitly as the sentinel mechanism's known limit, backstopped
outside this model by KV TTL and the 120-second presence-liveness window
--- neither of which is given a live transition here, in the same spirit
in which \texttt{nats-relay.tex}'s $DeleteSession$ leaves TTL expiry as
commentary rather than an operation.

\begin{schema}{CrashNoSentinel}
  \Delta SentinelWorld
  \where
  exited = zfalse \\
  own \in rowPresent \\
  %%
  exited' = ztrue \\
  crashed' = ztrue \\
  %%
  rowPresent' = rowPresent \\
  sentinelPresent' = sentinelPresent \\
  sentinelWriteAttempted' = sentinelWriteAttempted \\
  shutdownSignaled' = shutdownSignaled \\
  releaseAttempted' = releaseAttempted
\end{schema}

%%----------------------------------------------------------------------
\subsection{ResumeOwn}
%%----------------------------------------------------------------------

A fresh incarnation registers under the \emph{same}, resume-stable
$own$ key --- \texttt{register\_session}'s claim-then-write, whose
\texttt{update\_session} call overwrites the KV row unconditionally
regardless of what was there before.  This is one of the two
``accidental backstops'' the bead names; here it is encoded as the
actual overwrite the code performs, not assumed safe.

\begin{schema}{ResumeOwn}
  \Delta SentinelWorld
  \where
  exited = ztrue \\
  %%
  rowPresent' = rowPresent \cup \{ own \} \\
  sentinelWriteAttempted' = zfalse \\
  shutdownSignaled' = zfalse \\
  releaseAttempted' = zfalse \\
  exited' = zfalse \\
  %%
  sentinelPresent' = sentinelPresent \\
  crashed' = crashed
\end{schema}

Resume does not itself touch $sentinelPresent$: \texttt{register\_
session} claims the TTY name and overwrites the KV row, but never reads
or writes the sentinel directory.  A stale sentinel for $own$ left by
the departed prior incarnation is cleared later, by the resumed
incarnation's own first $OwnOrphanReapTick$ --- \texttt{\_active\_
lifespan} calls \texttt{\_reap\_sentinels} explicitly, immediately after
\texttt{register\_session} and before the periodic reaper task is even
created.  $ResumeOwn$'s postcondition $rowPresent' \ni own$ is what
makes that later discard harmless: see \S\ref{sec:resume-safe}.

%%----------------------------------------------------------------------
\section{Operations on Foreign Keys}
%%----------------------------------------------------------------------

The remaining three operations model a session key other than $own$
--- a different, already-departing incarnation, reaped correctly
because the reaping incarnation's key differs from the sentinel's key
(the bead's ``never-resumes case'').

%%----------------------------------------------------------------------
\subsection{RegisterForeign}
%%----------------------------------------------------------------------

\begin{schema}{RegisterForeign}
  \Delta SentinelWorld \\
  sk? : SESSIONKEY
  \where
  sk? \neq own \\
  sk? \notin rowPresent \\
  %%
  rowPresent' = rowPresent \cup \{ sk? \} \\
  %%
  sentinelPresent' = sentinelPresent \\
  sentinelWriteAttempted' = sentinelWriteAttempted \\
  shutdownSignaled' = shutdownSignaled \\
  releaseAttempted' = releaseAttempted \\
  exited' = exited \\
  crashed' = crashed
\end{schema}

%%----------------------------------------------------------------------
\subsection{WriteForeignSentinel}
%%----------------------------------------------------------------------

\begin{schema}{WriteForeignSentinel}
  \Delta SentinelWorld \\
  sk? : SESSIONKEY
  \where
  sk? \neq own \\
  %%
  sentinelPresent' = sentinelPresent \cup \{ sk? \} \\
  %%
  rowPresent' = rowPresent \\
  sentinelWriteAttempted' = sentinelWriteAttempted \\
  shutdownSignaled' = shutdownSignaled \\
  releaseAttempted' = releaseAttempted \\
  exited' = exited \\
  crashed' = crashed
\end{schema}

%%----------------------------------------------------------------------
\subsection{OwnReapForeign / OwnReapForeignFails}
%%----------------------------------------------------------------------

The good path: $own$'s reaper finds a sentinel for a key that is not
its own and reaps it fully --- \texttt{\_reap\_dead\_session} deleting
the row, releasing the TTY reservation, appending the logout event, and
the sentinel being removed on success.  $OwnReapForeignFails$ models
\texttt{delete\_session} raising: both row and sentinel are left in
place for a later retry, never producing an orphaned row.

\begin{schema}{OwnReapForeign}
  \Delta SentinelWorld \\
  sk? : SESSIONKEY
  \where
  sk? \neq own \\
  sk? \in sentinelPresent \\
  sk? \in rowPresent \\
  %%
  rowPresent' = rowPresent \setminus \{ sk? \} \\
  sentinelPresent' = sentinelPresent \setminus \{ sk? \} \\
  %%
  sentinelWriteAttempted' = sentinelWriteAttempted \\
  shutdownSignaled' = shutdownSignaled \\
  releaseAttempted' = releaseAttempted \\
  exited' = exited \\
  crashed' = crashed
\end{schema}

\begin{schema}{OwnReapForeignFails}
  \Delta SentinelWorld \\
  sk? : SESSIONKEY
  \where
  sk? \neq own \\
  sk? \in sentinelPresent \\
  %%
  rowPresent' = rowPresent \\
  sentinelPresent' = sentinelPresent \\
  %%
  sentinelWriteAttempted' = sentinelWriteAttempted \\
  shutdownSignaled' = shutdownSignaled \\
  releaseAttempted' = releaseAttempted \\
  exited' = exited \\
  crashed' = crashed
\end{schema}

%%----------------------------------------------------------------------
\section{Properties and Model Checking}
\label{sec:modelcheck}
%%----------------------------------------------------------------------

The property the bead asks for is

$$BadOrphan \;\widehat{=}\; own \in rowPresent \land own \notin sentinelPresent
\land exited = ztrue,$$

read as ``$own$'s row has no live process behind it and no sentinel
left to reclaim it.''  Two reachability goals are checked against it.

\subsection{(G1) The known gap: crash reaches $BadOrphan$ in one step}
\label{sec:crash-goal}

\begin{verbatim}
probcli docs/sentinel-reap.tex -model_check -noass -nodead \
  -card SESSIONKEY 2 \
  -goal "own : rowPresent & not(own : sentinelPresent) & exited=ztrue"
\end{verbatim}

ProB finds the goal at the shortest possible distance --- one
operation, $Init$; $CrashNoSentinel$.  $Init$ already has $own \in
rowPresent$ and
$exited = zfalse$, so the crash operation is enabled immediately.  This
is not a new finding --- it is the mechanism doing exactly what
biff-teh's own discovery notes already say it does: a process killed
without warning gets no sentinel, and the row survives until KV TTL or
the liveness window closes it, neither of which this specification
models as a live transition.  It is reported here because the bead asks
the literal property to be checked, not assumed, and because it fixes
the baseline against which \S\ref{sec:race} is the interesting result.

\subsection{(G2) The new finding: the fallback sentinel races its own reaper}
\label{sec:race}

Restricting the goal to $crashed = zfalse$ rules out (G1)'s trivial
witness and forces ProB to search for a route to $BadOrphan$ that goes
through no crash at all:

\begin{verbatim}
probcli docs/sentinel-reap.tex -model_check -noass -nodead \
  -card SESSIONKEY 2 \
  -goal "own : rowPresent & not(own : sentinelPresent) \
         & exited=ztrue & crashed=zfalse"
\end{verbatim}

ProB finds it, five operations from $Init$:

\begin{enumerate}
  \item $WriteOwnSentinel$ --- $own$'s incarnation begins an ordinary,
    non-signal shutdown; \texttt{\_write\_reap\_fallback\_sentinels}
    writes the sentinel.  $sentinelPresent = \{ own \}$,
    $sentinelWriteAttempted = ztrue$.
  \item $OwnOrphanReapTick$ --- before \texttt{\_shutdown\_tasks} has
    run, $own$'s \emph{own} background \texttt{\_reap\_loop} ticks
    (its two-second interval elapses during one of the \texttt{await}s
    in \texttt{\_append\_logout\_event} / \texttt{\_append\_companion\_
    logout\_event}, each bounded at three seconds), finds its own
    sentinel, and discards it under the own-key skip.
    $sentinelPresent = \emptyset$; $rowPresent$ is untouched.
  \item $SignalShutdown$ --- \texttt{\_shutdown\_tasks} now runs,
    stopping the reap loop for good.  Enabled because
    $sentinelWriteAttempted = ztrue$ from step 1, regardless of the
    sentinel's fate since.
  \item $ReleaseOwnFails$ --- \texttt{\_release\_session}'s
    \texttt{delete\_session} call times out or raises within its
    three-second bound.  $rowPresent$ still contains $own$;
    $sentinelPresent$ is still empty.
  \item $ExitProcess$ --- \texttt{\_lifespan\_cleanup} finishes anyway
    (every step in it is best-effort) and the process exits.
\end{enumerate}

The resulting state has $own \in rowPresent$, $own \notin
sentinelPresent$, $exited = ztrue$, $crashed = zfalse$: $BadOrphan$,
reached with no signal, no crash, and no resume anywhere in the trace.
This is a genuine defect internal to the mechanism \#350 ships, not a
restatement of the already-known SIGKILL gap: it is the interaction of
two of \#350's own changes --- an unconditional, unguarded fallback
sentinel write early in \texttt{\_lifespan\_cleanup}, and a per-step
teardown bound that lets \texttt{delete\_session} fail without
retrying inline --- with a reap loop that was never told to stand down
until \texttt{\_shutdown\_tasks} runs several awaits later.  Nothing in
the shipped code prevents $own$'s own periodic tick from firing in that
window; \texttt{\_reap\_loop}'s only guard is
\texttt{shutdown.is\_set()}, which is false for the entire duration of
steps 1--3.  \textbf{This is a finding for the operator to route, not a
defect this specification task fixes.}  The Python-level repair belongs
to whoever owns \texttt{app.py}: candidates include stopping the reap
loop (or at minimum excluding $own$'s own key from
\texttt{\_reap\_sentinels}) before \texttt{\_write\_reap\_fallback\_
sentinels} runs, or making the own-key skip conditional on
$shutdownSignaled$ rather than unconditional.

\subsection{Resume is safe by construction, not by the two named backstops}
\label{sec:resume-safe}

The bead poses a specific question: is the resume case ``correct by two
accidental backstops'' (\texttt{register\_session}'s stale-TTY
self-release and \texttt{update\_session}'s unconditional overwrite),
or is it correct for a more principled reason?  The model answers
directly, without assuming the answer.  $rowPresent$ is touched by
exactly two operations: $ResumeOwn$, which always \emph{adds} $own$,
and $ReleaseOwnSuccess$, which always \emph{removes} it --- and removes
$own$ from $sentinelPresent$ in the very same step.  So the only way to
reach a state with $own \in rowPresent \land own \notin sentinelPresent$
is for $ReleaseOwnSuccess$ to \emph{not} have been the most recent
sentinel-affecting event on $own$; in particular, immediately after
$ResumeOwn$, $exited' = zfalse$, so $BadOrphan$'s third conjunct is
false regardless of what $OwnOrphanReapTick$ does next.  A stale
sentinel discarded by the resumed incarnation's own reap tick is
harmless not because the TTY-reservation release or the KV overwrite
happen to compensate for it, but because $ResumeOwn$'s postcondition
already re-establishes $own \in rowPresent$ \emph{before} any reap tick
can run --- \texttt{\_active\_lifespan} calls \texttt{register\_session}
and then \texttt{\_reap\_sentinels} in that fixed order, and the
periodic reaper task does not even exist until afterwards.  Confirmed
by a third goal:

\begin{verbatim}
probcli docs/sentinel-reap.tex -model_check -noass -nodead \
  -card SESSIONKEY 2 \
  -goal "own : rowPresent & not(own : sentinelPresent) \
         & exited=ztrue & crashed=zfalse & shutdownSignaled=zfalse"
\end{verbatim}

This restricts the search to states reachable \emph{without} a
completed $SignalShutdown$ since the state last had $exited = ztrue$
--- i.e., states that could only have arisen from a resume path, since
$ResumeOwn$ is the only operation that clears $shutdownSignaled$ back
to $zfalse$ once $exited$ has been $ztrue$.  ProB visits the entire
reachable space at $\mathtt{-card\ SESSIONKEY\ 2}$ (136 states) and
reports no counterexample: the resume path never reaches $BadOrphan$.
The two mechanisms the bead named are real and both fire in production,
but neither is what makes resume safe in this model; what makes it safe
is the fixed call order in \texttt{\_active\_lifespan}, which the model
encodes as $ResumeOwn$'s postcondition running strictly before any
$OwnOrphanReapTick$ can matter.

\subsection{Foreign reaping never partially strands a row}

This one is read off the schemas rather than searched for, in the
manner of \texttt{session-model.tex}'s ``two facts~\ldots~here by
construction rather than as invariants to be checked'' (\S 5.2 there).
For a foreign key $sk?$, exactly one operation ever removes it from
$sentinelPresent$: $OwnReapForeign$, whose single effect clause removes
it from $rowPresent$ \emph{in the same schema}, i.e.\ the same atomic
state transition.  There is no foreign-key counterpart to
$OwnOrphanReapTick$ --- no operation discards a foreign sentinel while
leaving its row behind.  $OwnReapForeignFails$, the only other operation
touching a foreign sentinel, leaves \emph{both} $rowPresent$ and
$sentinelPresent$ untouched, so it can only prolong the recoverable
``row and sentinel both present, awaiting retry'' state, never produce
one where the sentinel is gone and the row is not.  The own-key race of
\S\ref{sec:race} is therefore specific to $own$: it exploits the split
between $OwnOrphanReapTick$ (sentinel-only) and $ReleaseOwn{}^*$
(row-and-sentinel, gated separately on $shutdownSignaled$), a split that
exists only because a reaper is never allowed to fully reap its own key
--- the very split the own-key skip introduces to protect the resume
case.  The foreign path has no such split and so no such race.

%%----------------------------------------------------------------------
\section{Model-Checking Scope}
%%----------------------------------------------------------------------

The full state space is checked exhaustively at $\mathtt{-card\
SESSIONKEY\ 2}$ --- one element for $own$, one for a foreign key,
enough to exercise every operation including $RegisterForeign$,
$WriteForeignSentinel$, $OwnReapForeign$, and
$OwnReapForeignFails$, which all need a second $SESSIONKEY$ distinct
from $own$:

\begin{verbatim}
probcli docs/sentinel-reap.tex -model_check -noass -nodead \
  -card SESSIONKEY 2 -p MAX_OPERATIONS 20000 -p TIME_OUT 60000 -coverage
\end{verbatim}

probcli visits $137$ states ($136$ live plus the search root) via $699$
fired transitions, reports \texttt{ALL OPERATIONS COVERED}, zero
deadlocks, and zero invariant violations (there being no state
invariant left to check --- \S\ref{sec:intro} explains why).  Every
operation is enabled from some reachable state and $ResumeOwn$ alone is
always enabled whenever $exited = ztrue$, so the system can never get
permanently stuck once a process has gone, crashed or not.

%%----------------------------------------------------------------------
\section{Discussion}
%%----------------------------------------------------------------------

\subsection{What this model captures}

\begin{itemize}
  \item Sentinel writing, unified across the signal-handler path and
    the \#350 fallback path as one operation ($WriteOwnSentinel$),
    matching the code's own convergence onto \texttt{\_write\_sentinel}.
  \item Sentinel consumption, split into the own-key skip
    ($OwnOrphanReapTick$, no row effect) and the foreign path
    ($OwnReapForeign$/$OwnReapForeignFails$, atomic row-and-sentinel
    removal on success).
  \item KV row deletion ($ReleaseOwnSuccess$, $OwnReapForeign$) and
    row orphaning --- the state $BadOrphan$ that both goals target.
  \item The liveness relationship between a sentinel and the row it
    refers to: a sentinel's only job is to trigger the row's deletion;
    \S\ref{sec:race} shows that job can be discharged (the sentinel
    consumed) without the row actually being deleted, which is the
    exact failure mode the mechanism exists to prevent.
\end{itemize}

\subsection{What this model does not cover}

\begin{itemize}
  \item \textbf{KV TTL and the presence-liveness window} --- the
    backstop for $CrashNoSentinel$, deliberately left as commentary
    rather than a live transition, following \texttt{nats-relay.tex}'s
    own treatment of $DeleteSession$.
  \item \textbf{Multiple concurrent reaping servers} --- the model
    takes the perspective of one incarnation reaping foreign keys; two
    servers racing to reap the \emph{same} foreign sentinel is a
    different, narrower race (both attempt $OwnReapForeign$ on the same
    $sk?$; the loser's \texttt{delete\_session} finds nothing to
    delete) that the code's own commentary in \texttt{\_handle\_kv\_
    delete} already accepts as harmless duplication, not a correctness
    hazard for this specification's property.
  \item \textbf{Companion sessions as a distinct carrier} --- folded
    into $own$ by the remark in \S\ref{sec:own}, not given a second
    constant; every property proved of $own$ transfers unchanged.
  \item \textbf{The content of a KV row} --- $rowPresent$ is presence
    only, matching this document's use of \texttt{nats-relay.tex}'s
    $SessionStore$ for the underlying storage model.
\end{itemize}

\end{document}
